% =============================================================================
% PreCex：反例驱动的综合前 RTL 缺陷定位与修复
% 《计算机辅助设计与图形学学报》投稿格式
% 编译：LuaLaTeX（scripts/compile.ps1 -Name precex_zh）
% 源稿：02_Precex/paper/manuscript/precex_paper.tex（含历轮审稿修订）
% =============================================================================
\documentclass[twocolumn,10pt]{ctexart}

% ---------- 字体（项目内 fonts/；compile.ps1 设 OSFONTDIR） ----------
\setCJKmainfont{FZShuSong-Z01}       % 方正书宋_GBK（中文正文）
\setCJKsansfont{SimHei}              % 黑体
\newCJKfontfamily\songfamily{FZShuSong-Z01}
\newCJKfontfamily\heifamily{SimHei}
\newCJKfontfamily\fangfamily{FZFangSong-Z02}  % 方正仿宋_GBK
\let\songti\songfamily
\let\heiti\heifamily
\let\fangsong\fangfamily
\setmainfont{Times New Roman}        % 英文/数字

% ---------- 宏包 ----------
\usepackage{amsmath,amssymb}
\usepackage{newtxmath}
\usepackage{booktabs}
\usepackage{graphicx}
\usepackage{array}
\usepackage{cite}
\usepackage[hidelinks]{hyperref}
\usepackage{xurl}
\usepackage{geometry}
\usepackage{caption}
\usepackage{subfig}

% 版面：A4 双栏
\geometry{a4paper,left=2.0cm,right=2.0cm,top=2.5cm,bottom=2.5cm,columnsep=31.5pt}

% 标题层级（1级11.5pt黑体 / 2级五号黑体 / 3级10pt书宋）
\ctexset{
  section = {format=\heiti\fontsize{11.5}{15}\selectfont, beforeskip=6pt, afterskip=4pt},
  subsection = {format=\heiti\fontsize{10.5}{14}\selectfont, beforeskip=5pt, afterskip=3pt},
  subsubsection = {format=\songti\fontsize{10}{13}\selectfont, beforeskip=4pt, afterskip=2pt}
}

% 图题（书宋9pt）、表题（黑体9pt）、表注/分图题（书宋8pt）
\DeclareCaptionFont{songti}{\songti}
\DeclareCaptionFont{heiti}{\heiti}
\captionsetup[figure]{font={songti,small},labelfont={bf}}
\captionsetup[table]{font={heiti,small},labelfont={bf}}
\captionsetup[subfloat]{font={songti,scriptsize}}
\captionsetup{labelsep=space}

\graphicspath{{../figures/final/}}
\setlength{\emergencystretch}{3em}
\Urlmuskip=0mu plus 3mu\relax

\begin{document}

\pagestyle{plain}

% ---------------------------------------------------------------------------
% 中文题录（通栏）
% ---------------------------------------------------------------------------
\twocolumn[%
\begin{center}
  {\heiti\zihao{3} PreCex：反例驱动的综合前 RTL 缺陷定位与修复}\par
  \vspace{6pt}
  {\fangsong\zihao{4} 佟亚龙$^{1)}$}\par
  \vspace{4pt}
  {\songti\fontsize{8}{12}\selectfont（国防科技大学计算机学院，长沙 410073）}\par
  \vspace{2pt}
  {\songti\fontsize{8}{12}\selectfont E-mail：tyl2003@nudt.edu.cn（通信作者）}\par
  \vspace{8pt}
\end{center}
\begin{flushleft}
  {\heiti\fontsize{8.5}{13}\selectfont 摘\quad 要：}{\songti\fontsize{8.5}{13}\selectfont
  针对弱测试平台仿真通过而形式验证失败的跨周期寄存器传输级（RTL）缺陷在综合前难以定位与
  修复的问题，提出 PreCex——一种完全开源、反例驱动的 RTL 缺陷定位与修复流水线
  （iverilog/Yosys/SymbiYosys/Z3 加大语言模型，LLM）。所提方法将失败的形式化反例转换为
  结构化证据，定位缺陷区域，生成最小补丁，并在 golden-first 的有界模型检验（BMC）下重新
  验证。在覆盖 6 个模块、7 类错误的 34 样本 L3 基准上，以 DeepSeek v4-flash 完成无泄漏
  四设置 408 次受控实验：统一 BMC 判据下四种证据表征的回归通过率均达 100\%；卫生化前
  定位精度为 A 87.3\%、B 75.5\%、D 73.5\%、C 60.8\%（102/102，行号判据修复后口径），
  A 与 C、A 与 D 差异显著（Holm 校正后 p=4.5e-5、p=0.047）；仅修正行号基准与断言行语义
  等呈现问题后，四设置定位精度收窄至 A 95.1\%、D 95.1\%、B 94.1\%、C 92.2\%（极差由
  29.4~pp 降至 2.9~pp），配对检验全部不显著。卫生化后成本 D 最低（0.24 美元/102 次）、
  A 最高（0.37 美元）。变异分析杀死 88.5\% 强变异体与 81.8\% 常量变异体，独立 T2 审计
  408/408 通过。进一步报告评测卫生教训：三处 ground-truth 泄漏曾致“B 显著最优”伪结论
  （剥离后 134 处定位翻转），prove/k-归纳判据误判 78 个正确修复，故采用 golden-first 的
  BMC 判据；L2 门禁对照（72 次）回归通过率 91.7\%，弱 tb 门禁不构成能力分界。}\par
  \vspace{4pt}
  {\heiti\fontsize{8.5}{13}\selectfont 关键词：}{\songti\fontsize{8.5}{13}\selectfont
  RTL 调试；形式验证；反例；大语言模型修复；变异分析}\par
  \vspace{2pt}
  {\songti\fontsize{8.5}{13}\selectfont 中图法分类号：TN47\qquad DOI：10.3724/SP.J.1089.2026-00001}\par
  \vspace{6pt}
\end{flushleft}

% ---------------------------------------------------------------------------
% 英文题录
% ---------------------------------------------------------------------------
\begin{center}
  {\bfseries\fontsize{13.5}{18}\selectfont PreCex: Counterexample-Driven Pre-Synthesis RTL Defect Localization and Repair}\par
  \vspace{5pt}
  {\normalsize Yalong Tong$^{1)}$}\par
  \vspace{3pt}
  {\small\itshape (College of Computer Science and Technology, National University of Defense Technology, Changsha 410073, China)}\par
  \vspace{6pt}
\end{center}
\begin{flushleft}
  {\bfseries Abstract:} {\small Cross-cycle behavioral defects -- RTL bugs that pass weak testbench
  simulation yet fail formal verification -- are among the hardest problems to localize and repair
  before synthesis. This paper presents PreCex, a fully open-source, counterexample-driven RTL bug
  localization and repair pipeline (iverilog/Yosys/SymbiYosys/Z3 with an LLM). It converts failing
  formal counterexamples into structured evidence, localizes the defective region, generates a minimal
  patch, and re-verifies under bounded model checking (BMC) with a golden-first rule. On a 34-sample
  L3 benchmark covering 6 modules and 7 error classes, a controlled four-setting three-seed ablation
  (408 leak-free DeepSeek v4-flash runs) answers how evidence representation affects localization and
  repair, and how evaluation reliability is determined by data and criterion hygiene. Under a unified
  BMC criterion, all four representations reach a 100\% regression pass rate. Before hygiene,
  localization accuracy is A 87.3\%, B 75.5\%, D 73.5\%, C 60.8\% (102/102, line-number-criterion-
  fixed basis), with A significantly higher than C (McNemar p=7.4e-6, Holm p=4.5e-5) and D
  (Holm p=0.047); after a four-setting hygiene ablation fixing only presentation issues (line-number
  basis and assertion-line semantics), accuracy narrows to A 95.1\%, D 95.1\%, B 94.1\%, C 92.2\%
  (range from 29.4~pp to 2.9~pp), all paired tests non-significant. Hygiene-cost per 102 runs is
  D 0.24, A 0.37 USD. Mutation analysis kills 88.5\% of strong and 81.8\% of constant mutants; an
  independent T2 audit passes 408/408. A data-hygiene lesson is reported: three ground-truth leaks
  produced a false “B significantly best” conclusion (134 outcomes flipped after removal), and
  prove/k-induction misjudges 78 correct repairs on gated-temporal-assertion modules, motivating
  the golden-first BMC criterion. An L2 gate control (72 runs) reports a 91.7\% regression pass
  rate; the weak-testbench gate is not a capability boundary.}\par
  \vspace{3pt}
  {\bfseries Key words:} {\small RTL debugging; formal verification; counterexample; LLM-based repair; mutation analysis}\par
\end{flushleft}
]%

% ---------------------------------------------------------------------------
% 正文（双栏）
% ---------------------------------------------------------------------------

% ---------- 引言（不加章节号与标题） ----------
\subsection*{背景与问题}

寄存器传输级（RTL）设计复杂度持续增长，跨周期行为缺陷（L3：弱测试平台通过而形式验证失败）
是综合前最难定位与修复的缺陷类型之一。此类缺陷通常隐藏在状态机迁移、握手时序或边界回绕
逻辑中：简单断言行启发式无法命中注入点，仿真波形中的信息过载使人工与自动化定位都极为耗时。
一旦缺陷进入综合与流片流程，修复成本呈数量级上升，因此在综合前准确、低成本地定位并修复
RTL 缺陷具有明确工程价值。

需要精确界定本文的适用边界：PreCex 主要针对失效轨迹较短（本基准中位 6 拍）、可通过局部
代码修正恢复的 L3 缺陷子集；对于长周期状态机死锁、多周期握手超时等需要结构性重写的缺陷，
本系统生成的候选补丁虽能通过有限深度 BMC 回归，但属于工程上的绕过方案，完整修复留待未来
工作（见 \S\ref{sec:limitation}）。

现有基于大语言模型（LLM）的修复流水线主要面临两类困难。其一，直接输入原始日志或波形，
信息过载导致定位噪声，模型难以从数千行仿真记录中筛出真正的因果信号。其二，依赖参考模型
（golden reference model）进行差分调试的方法在综合前阶段往往不可用——此时不存在可作为
规范的金牌实现。近期工作（FVDebug、GROVE）证明了结构化反例理解的价值，但仍缺乏四点能力：
（1）开源端到端流水线；（2）面向跨周期缺陷的系统性研究；（3）严格的验证充分性闭环；
（4）证据表征的受控消融。

\subsection*{设计动机}

本文观察到，传递给 LLM 的\textbf{证据表征}（evidence representation）是一个关键但缺乏研究
的维度：同样的模型能力下，证据如何组织直接决定定位精度与 token 成本。为此，我们在同一
34 样本 L3 数据集上，用 DeepSeek v4-flash 对三种证据表征做三-seed 严格配对的受控消融
（A/B/C，306 次无泄漏运行（C 102/102），见 \S\ref{sec:setup}）：A 为原始反例日志与 VCD
头尾；B 为结构化证据（EvidenceEngine JSON）；C 为反例语义化（CexSemantizer 输出的周期事件、
状态轨迹、故障锥与自然语言摘要）；D 为静态故障锥（static fault cone）（确定性提取），作为
探索性设置复跑。初版实验曾因三处 ground-truth 泄漏（\S\ref{sec:hygiene}）得到“B 显著最优”
的伪结论，剥离泄漏后结论反转，这本身成为本文的评测卫生教训。

与此同时，本文强调修复判据本身必须接受检验：原协议采用 prove/k-归纳判据，但对带门控时序
断言的 axi/uart 模块不收敛——golden 设计本身即返回 UNKNOWN，导致 78 个正确修复被误判为
失败。据此我们确立了\textbf{golden-first 规则}：任何修复判据必须首先在 golden 设计上通过，
才可用于评判修复结果。

\subsection*{贡献}

本文的主要贡献如下：

\begin{itemize}
\item \textbf{开源反例驱动修复流水线与评测资产}：基于 iverilog/Yosys/SymbiYosys/Z3 与
DeepSeek v4-flash 的端到端实现，完全可复现；配套 34 样本 L3 数据集，含 golden 双重校验、
ground-truth 隔离审计与充分性审计（变异/空泛/T2/delta-debugging）。
\item \textbf{证据呈现卫生主导定位精度（全四设置卫生消融）}：统一 BMC 判据下四种设置在
实验筛选口径下的 BMC 回归通过率均达 100\%；卫生化前定位精度（A 87.3\%、B 75.5\%、D 73.5\%、
C 60.8\%，行号判据修复后口径）呈现 A 与 C、A 与 D 的显著差异，但仅修正行号基准与断言行
语义等呈现问题（全四设置 408 次）后，定位精度收窄为 A 95.1\%、D 95.1\%、B 94.1\%、C 92.2\%
（极差从 29.4~pp 降至 2.9~pp），配对检验全部不显著（A/C p=0.375、A/D p=1.0）——证据呈现
卫生而非证据表征本身主导了卫生化前的差异；成本（卫生化口径）D 最低（\$0.24），B 次之
（\$0.29），A 最高（\$0.37）。
\item \textbf{评测双卫生方法论教训}：ground-truth 泄漏（三处）会使“B 显著最优”成为伪结论
（剥离后 134 处定位翻转、结论反转）；prove/k-归纳判据失效会误判 78 个正确修复。数据卫生
（ground-truth 隔离）与判据卫生（golden-first）是自动修复评测的可靠基准。
\end{itemize}

% ---------- 1 级标题 ----------
\section{相关工作}

\subsection{反例驱动的调试与根因分析}

\textbf{FVDebug}（Bai 等，2025）\cite{fvdebug}是最接近的同类工作：它从失败的形式化轨迹构建
因果 DAG，用批量 LLM 调用扫描可疑节点，通过 agent 生成根因叙述并输出修复，在两个工业反例上
完成演示。PreCex 与其共享反例理解这一核心：我们的设置 D 是静态故障锥（static fault cone）的
确定性开源实现（失败断言、故障锥、全周期状态轨迹与触发条件）。FVDebug 在三个维度上留有空缺，
而 PreCex 恰好控制：其一，它依赖商业工具链；其二，它没有报告证据表征的受控消融；其三，它不
量化验证器充分性。PreCex 完全开源（iverilog/Yosys/SymbiYosys/Z3），在同一 34 样本数据集上
比较四种证据表征，并以变异分析、空泛性控制与独立 T2 审计评测验证器。软件缺陷定位（fault
localization）是硬件与软件调试的共同基础问题，中文期刊已有大量研究：王浩仁等基于冗余覆盖
信息约简提升可疑度排序\cite{zh1}，曹鹤玲与姜淑娟以 Chameleon 聚类处理多错误定位\cite{zh2}，
王建峰等利用错误交互集完成组合测试故障定位\cite{zh3}；这些工作聚焦软件频谱信息，本文将其
方法论迁移到 RTL 反例证据组织，并额外引入形式验证判据闭环。

\textbf{GROVE}（Bai 与 Ren，2025）\cite{grove}将断言失败调试知识组织为分层知识树：训练时以
模型检验验证知识项，测试时按预算感知检索，一致提升 pass@1/pass@5。GROVE 与 PreCex 互补：
GROVE 跨样本积累可复用知识，PreCex 消费单个失败反例，在无历史存储的闭环中完成定位-修复-重验。
GROVE 式知识层可作为本流水线的自然扩展。

\textbf{开源 LLM 驱动形式验证}（Tran，2026）\cite{openllmformal}采用与本文相同的开源工具链
（Yosys/SymbiYosys/Z3），以 LLM 引导反例驱动迭代修复，以 k-归纳证明或预算耗尽为停止条件，
在 6 个基准上仅可靠修复 1 个。该工作验证了开源管线的可行性并整理了失败模式；PreCex 的差异
在于增加证据语义化层（该工作将原始反例直接送入 LLM），并在更大、受控的 L3 数据集上采用
golden-first 判据。

\textbf{FormalRTL}\cite{formalrtl}以软件参考模型作为形式化规范，闭合规划、综合与等价检验循环，
并含一个反例简化器，只高亮失败功能与不匹配信号。PreCex 不作参考模型假设：仅凭断言与反例工作，
这正是综合前无黄金模型时的常见设置。

\textbf{ChipAgents RCA}（商业，2026）\cite{chipagents}将波形理解引擎与 prover-verifier agent
对及自一致性排序结合，报告了工业部署中对三周期竞争条件的定位。它佐证跨周期根因分析是真实
工业痛点，但作为闭源黑盒，未提供任何消融细节。

\subsection{基于 LLM 的 RTL 修复}

\textbf{VeriPilot}\cite{veripilot}使用黄金参考模型与 CDFG 信号追踪，将 CVDP 上的 GPT-4o
修复率从 54.3\% 提升至 85.71\%。\textbf{VeriDebug}\cite{veridebug}学习对比嵌入并引导修正，
在定位加类型联合任务上取得 64.7\% 的 Acc@1。\textbf{RTLFixer}\cite{rtlfixer}采用检索增强
修复语法错误。三者均非反例驱动：分别依赖参考模型、微调嵌入或仅做语法级修复。PreCex 明确
针对综合前无参考模型的场景，将证据表征而非模型能力作为研究变量。国内基于深度学习与频谱
信息的缺陷定位方法亦有系统研究\cite{zh5}，本文聚焦 RTL 反例驱动场景，与软件侧方法形成互补。

\textbf{Fixbench-RTL}\cite{fixbench}与\textbf{HDL-FixBench}\cite{hdlfixbench}是修复基准：
Fixbench-RTL 覆盖语法、功能与安全域；HDL-FixBench 面向仓库级实例（OpenTitan/CVA6/Ibex 共
57 个实例，最佳模型 40.3\%）。HDL-FixBench 所暴露的缺陷正是 BugBench-PS 着力规避的：规模
过小、多样性不足、补丁阈值存在争议、缺少非 LLM 基线以及数据污染风险。

\subsection{断言生成与验证充分性}

\textbf{AssertLLM}\cite{assertllm}与\textbf{AssertLLM2}\cite{assertllm2}从自然语言规格加波形
生成 SVA 断言；AssertLLM2 引入 83 设计的断言基准，并以变异驱动的 bug-hunting 设置评测断言。
\textbf{AssertGen}\cite{assertgen}通过思维链提取验证目标并桥接跨层信号，采用变异测试覆盖率
评估。\textbf{LASSO}\cite{lasso}生成带显式空泛性检查与覆盖反馈的安全属性，在 OpenTitan 中
发现 5 个真实缺陷。\textbf{LintLLM}\cite{lintllm}以变异缺陷注入执行 LLM 代码检查。

这些工作生成或审计断言；PreCex 消费断言加反例以完成定位与修复。缺陷定位评测的可信度同样
取决于数据与判据的卫生性，中文软件工程界对此已有丰富实践\cite{zh4}。其评测实践（变异覆盖、
空泛性、bug-hunting）构成了本文充分性审计的方法学先例：88.5\% 强变异与 81.8\% 常量变异被杀，
删除变异空泛性控制通过。

\subsection{基准与评测方法}

工具执行裁决（tool-execution adjudication）是 EDA-AI 领域公认的评测实践：Rule2DRC\cite{rule2drc}
（ICML 2026）比较 13,921 条布局规则结果，PostEDA-Bench\cite{posteda}评估 145 个 DRC/PPA 任务，
均通过执行设计工具裁决。PreCex 遵循同一原则：修复成功由编译、弱测试平台回归与形式 BMC 裁决，
从不依赖 LLM 自述。工程语义数据集（OpenRTLSet\cite{openrtlset}、EDA-Schema-V2\cite{edaschema}、
ChipLingo\cite{chiplingo}）为本文证据模式设计提供参考。

\subsection{本文定位}

\begin{table}[!t]
\centering
\caption{设计空间定位：与已有工作的证据表征、开源性与评测闭环对比}
\label{tab:positioning}
\scriptsize
\setlength{\tabcolsep}{2pt}
\renewcommand{\arraystretch}{1.12}
\begin{tabular}{p{0.19\columnwidth}p{0.25\columnwidth}c p{0.13\columnwidth}p{0.11\columnwidth}p{0.13\columnwidth}}
\toprule
\textbf{工作} & \textbf{证据表征} & \textbf{开源} & \textbf{跨周期} & \textbf{消融} & \textbf{充分性} \\
\midrule
FVDebug & 故障锥 & 否 & 部分 & 否 & 否 \\
GROVE & 知识树 & 否 & 否 & 否 & 否 \\
Tran 2026 & 原始反例 & 是 & 部分 & 否 & 否 \\
FormalRTL & 反例简化+参考模型 & 部分 & 否 & 否 & 否 \\
VeriPilot & CDFG+参考模型 & 否 & 否 & 否 & 否 \\
VeriDebug & 学习嵌入 & 否 & 否 & 否 & 否 \\
AssertLLM2 & 断言生成 & --- & --- & --- & 变异驱动 \\
\textbf{PreCex} & \textbf{原始/结构化/语义化/因果图} & \textbf{是} & \textbf{34 样本} & \textbf{A/B/C/D} & \textbf{变异/空泛/T2} \\
\bottomrule
\end{tabular}
\end{table}

表~\ref{tab:positioning} 汇总了与已有工作的定位差异。PreCex 的三项差异化能力均有受控实验
佐证：(1) 完全开源、可复现的反例驱动流水线；(2) 四设置证据表征消融（卫生化前配对显著性 +
卫生化后收窄至无显著差异的完整报告）；(3) 带验证充分性闭环（变异、空泛、独立 T2 审计）的
跨周期 L3 基准。此外，判据翻案审计（第~\ref{sec:judgment}~节）贡献了一条方法论经验：修复
判据必须先用 golden 设计验证，才能用于评判修复。

\textbf{与反例精简文献的关系：}D 设置的静态故障锥（扇入并集 + 代码切片）在概念上接近
cone-of-influence（COI）精简、minimal unsatisfiable core（MUS）与反例最小化（counterexample
minimization）等经典形式化技术；与这些技术的差异在于：本文的静态扇入是\textbf{零 LLM 的
确定性启发式}，服务于\textbf{LLM 提示输入}而非求解器内核，且不提供完备性保证
（\S\ref{sec:limitation} 已声明）。该定位避免将静态扇入误解为新的形式化精简算法。

\textbf{跨领域启示：}本文的两条评测卫生教训并不局限于 RTL 修复——ground-truth 泄漏与判据-
黄金一致性检查在自动程序修复（APR）、代码生成评测与 LLM Agent 评测中均有对应物：任何以
参考实现或隐藏测试为监督的评测，都需要先验证判据本身在黄金参考上成立，并审计提示/证据是否
泄漏了求解目标。这使本文的方法论产出具备跨领域可迁移性。

% ---------- 方法 ----------
\section{方法}\label{sec:method}

\textbf{系统版本口径：}第~\ref{sec:method}~节描述的完整系统包含若干增量机制（历史修复记忆、
delta-debugging、自适应窗口与工具链预检），其中部分机制已实现但尚未完成全量重跑；主实验
（第~\ref{sec:setup}~节与第~\ref{sec:results}~节）基于一个功能子集（无历史修复记忆、固定
窗口 8、无 delta-debugging），该子集足以支撑本文的证据表征消融研究。为消除协议歧义，表~
\ref{tab:protocol} 给出各机制在不同实验阶段的启用清单。增量机制与主实验结论的兼容性由深
时序子集小规模实测佐证（第~\ref{sec:deep}~节，s38--s42 四设置 60/60 全部完成），完整系统的
更大规模全量评测留待后续工作。

\begin{table}[!t]
\centering
\caption{机制启用清单（主实验 / 卫生消融 / 深子集 / L2 对照）}\label{tab:protocol}
\scriptsize
\setlength{\tabcolsep}{1.5pt}
\renewcommand{\arraystretch}{1.15}
\begin{tabular}{p{0.24\columnwidth}p{0.18\columnwidth}p{0.18\columnwidth}p{0.16\columnwidth}p{0.16\columnwidth}}
\toprule
\textbf{机制} & \textbf{主实验} & \textbf{卫生消融} & \textbf{深子集} & \textbf{L2 对照} \\
\midrule
历史修复记忆（反馈循环） & 未启用 & 未启用 & 未启用 & 未启用 \\
自适应窗口（window\_eff） & 未启用（固定 window=8） & 未启用（固定 window=8） & \textbf{启用} & 未启用（固定 window=8） \\
delta-debugging（最小化） & 未启用（报告 LLM 直接输出） & 未启用 & 仅验证时参考 & 未启用 \\
行号卫生（绝对行号前缀等） & 未启用 & \textbf{启用} & 未启用 & 未启用 \\
工具链兼容性预检 & 启用 & 启用 & 启用 & 启用 \\
\bottomrule
\end{tabular}
\end{table}

\subsection{系统架构}

\begin{figure}[!htbp]
\centering
\includegraphics[width=0.98\columnwidth]{fig_pipeline.pdf}%
\caption{PreCex 流水线：失败的形式化反例经解析生成结构化证据（B），可选语义化（C）或汇总为
故障锥（D），由 LLM 修复器生成最小 diff，并在 golden-first 的 BMC 判据下重新验证；失败结果
回灌循环。}
\label{fig:pipeline}
\end{figure}

整体架构如图~\ref{fig:pipeline} 所示。PreCex 由四个核心组件与两个支撑元素构成（开源工具链
iverilog、Yosys\cite{yosys} 与 SymbiYosys\cite{sby}）。\textbf{EvidenceEngine} 将 cex.log
解析为统一 JSON 结构（错误类型、模块、文件、行号、代码片段、信号、触发条件、失败阶段与失败步）。
\textbf{CexSemantizer} 将 VCD 转换为周期事件表、状态轨迹、故障锥与自然语言摘要（window=8
压缩）。\textbf{LocalRepairer} 以设计、断言与证据为输入（DeepSeek v4-flash，temperature 0.2），
输出定位与最小统一 diff，重试不超过两轮；\textbf{每次失败尝试的 diff 与失败原因被注入下一轮
prompt}（历史修复记忆，见 \S\ref{sec:feedback}），明确要求 LLM 分析上次为何失败并避免重复该
模式，避免“开环重试”重复同一错误补丁。\textbf{Verifier} 执行三关验证：iverilog 编译零错误、
弱测试平台仿真全绿、sby smtbmc+z3 BMC 无反例（主判据），并叠加 golden 双重校验；工具链适配
由断言安全子集（Gate-1 收敛）与前置兼容性预检（scripts/check\_rtl\_compat.py，见
\S\ref{sec:toolchain}）保障，工具解析失败按 INCONCLUSIVE 记录而非崩溃。支撑元素为实验变量
控制组件与 BugBench-PS 数据集资产。

\subsection{证据引擎与反例语义化}\label{sec:evidence}

EvidenceEngine 以失败断言、错误类型与失败阶段的统一模式组织反例，输出稳定可解析的
evidence.json。CexSemantizer 在此基础上生成周期事件表（每周期信号变化）、状态轨迹（跨周期
信号演化）、故障锥与自然语言摘要。\textbf{故障锥的算法为静态代码级扇入}：取失败断言表达式
（trigger\_condition）、失败行 $\pm 4$ 行代码切片（code\_slice）与关键控制/状态/协议信号
（按模块时钟域枚举的 KEY\_SIGS）三者并集，再过滤参数/内部辅助前缀（DATA\_/OP\_/f\_ 等）。
它不依赖 VCD 中信号是否翻转——恒定错误值的根因信号仍经断言表达式/代码切片进入锥——但代价
是可能包含与缺陷无关的周边信号，该边界在 \S\ref{sec:limitation} 如实声明。实验表明主数据集
反例普遍较短（VCD 时钟事件合计 295 拍，中位 6 拍，最长 41 拍），window=8 窗口压缩的收益
有限——语义化的成本由摘要文本主导而非轨迹长度，这一观察直接支撑了第~\ref{sec:cost}~节的
成本分析。需要指出，固定 window=8 会截断长反例的关键因果尾部：最长样本 41 拍，而 8 拍之后
的窗口在状态机死锁/握手超时类缺陷上可能正是根因所在，属系统性信息损失而非成本权衡。为此
CexSemantizer 提供\textbf{失败时刻尾部回溯的自适应窗口}（window\_eff
$=\max(8,\min(24,\lfloor fail\_step/2\rfloor))$），深时序子集（35--75 拍，仓库 samples/deep）
的语义化已采用该策略；主数据集 34 样本为口径一致仍保留 window=8，该截断作为系统性局限列入
\S\ref{sec:limitation}。

需要指出：故障锥的静态扇入算法可能使 B 与 D 的证据包含较多与缺陷弱相关的周边信号——相比
之下，A（原始日志）不含故障锥，其信息密度虽低但冗余噪声也更少，这可能部分解释了主数据集上
A 的定位精度（87.3\%）高于 B（75.5\%）和 D（73.5\%）的现象。该假设的正确性需通过动态切片
或形式化轨迹最小化的消融实验验证，留待未来工作。

\subsection{四种证据表征}

\begin{table}[!t]
\centering
\caption{四种证据表征设置（受控消融）}
\label{tab:evidence}
\scriptsize
\setlength{\tabcolsep}{3pt}
\renewcommand{\arraystretch}{1.15}
\begin{tabular}{p{0.13\columnwidth}p{0.33\columnwidth}p{0.22\columnwidth}p{0.17\columnwidth}}
\toprule
\textbf{设置} & \textbf{证据} & \textbf{提取方式} & \textbf{单次成本} \\
\midrule
A & 原始 cex.log + VCD 头尾 & 无 & 基线 \\
B & evidence.json（结构化） & 确定性 & 低 \\
C & semantics.json（周期事件+状态轨迹+故障锥+摘要，window=8） & LLM 语义化 & 高 \\
D & 静态故障锥（static fault cone）（失败断言+故障锥+全周期状态轨迹+触发条件） & 确定性（零 LLM） & 最低 \\
\bottomrule
\end{tabular}
\end{table}

各设置的证据内容与提取成本见表~\ref{tab:evidence}。四种设置使用同一 prompt 家族，仅替换
证据段，以控制偏置。A 为基线，B 与 D 均为零 LLM 成本的确定性提取，C 引入一次 LLM 语义化
调用，是四者中唯一的额外推理成本来源。

\subsection{修复判据：golden-first 规则}

主判据为 BMC（verify.sby），与 golden 验证（verify\_golden.sby）一致。原协议采用 prove/k-
归纳（verify\_repair.sby）作为修复判定，但在带门控时序断言的 axi/uart 模块上不收敛：golden
设计本身返回 UNKNOWN（rc=4），旧判据下 78 个正确修复被误判为失败。同一正确修复（s33 的 B
设置 seed0、s17 的 B 设置 seed1）在 bmc 下 PASS、prove 下 FAIL 的实测直接暴露了这一矛盾。

据此确立 golden-first 规则：任何修复判据必须首先在 golden 设计上通过（判据-黄金一致性检查），
方可用于评判修复。对 A/B/C/D 408 次运行按 BMC 判据重验，408/408 全部通过（A/B/C 306 次中
303 条含 diff 记录直接重验、3 条网络失败补跑，D 102 条同协议复跑），旧 prove 判据下修复
通过率仅 73.5\%（225/306），修正为 BMC 判据下的 100\%。prove/k-归纳仅作为辅助参考保留，
不再作为失败判据。

\subsection{验证充分性审计}

为量化断言集能否捕获注入缺陷（充分性），本文对 golden 设计执行三类变异：
\begin{itemize}
\item \textbf{强变异}：比较器反转等强语义变异（d16），400 个变异体，88.5\% 被杀。
\item \textbf{常量变异}：替换参数常量 DATA\_W/ADDR\_W/DEPTH/DIV/HALF，484 个变异体，
81.8\% 被杀。
\item \textbf{删除变异}：移除断言，作为空泛性控制，全部通过（若删除断言后仍能杀变异，说明
断言冗余）。
\end{itemize}

模块级分析显示 axi 与 fifo 的强变异杀死率为 100\%，而 uart\_rx 最弱（强变异 55.0\%、常量
变异 40.9\%）——结构断言对波特时序常量漂移不敏感，这是已记录的局限并列为未来工作。此外，
T2 审计以确定性脚本替代 LLM 自证，对全部含 diff 的修复（A/B/C/D 共 408 条）检查三类违规：
接口/端口/参数变更、断言篡改、证据闭环一致性，均零失败；loc\_dev 中位数为 0（干净口径
A/B/C 74.5\% 精确命中），补丁规模均值 2.26 行、最大值 12 行（unified diff 增删行合计，
408 条）。\textbf{补丁“最小性”并非仅由规模数字佐证}：新增后验 delta-debugging 检查
（scripts/minimize\_patch.py，见 \S\ref{sec:minimality}）对通过修复迭代移除 diff 中的改动单元，
仅当移除后仍通过编译/仿真/BMC 三通过才保留，直至 1-minimal；在深时序子集真实修复上实测
确认无冗余单元（移除任一单元即 FAIL，见 \S\ref{sec:minimality}）。

\subsection{反馈循环与历史修复记忆}\label{sec:feedback}

\textbf{（未在主实验启用）}系统代码已实现反馈循环机制：每次失败尝试的 diff 与失败原因被写入
历史记录并注入下一轮 prompt，要求 LLM 分析上次失败后再给出不同方案。该机制面向后续扩展评估；
干净口径主实验（A/B/C/D 408 次）与 L2 对照（72 次）均运行于固定 prompt 协议（无历史注入），
报告的定位与修复数字反映该协议下的行为。

\subsection{最小补丁后验验证（delta-debugging）}\label{sec:minimality}

\textbf{（未在主实验启用）}为消除“最小”修补仅靠 LLM 提示词承诺的风险，实现了确定性
delta-debugging 组件（scripts/minimize\_patch.py）：将修补 diff 拆为改动单元，贪心剔除并
逐次 BMC 重验，保留的单元构成真实的 1-minimal 补丁。s40（fifo\_sync 边界回绕）经最小化后
patch 从 6 行缩至 4 行，独立通过 BMC 校验。主实验报告的修补为未最小化的 LLM 直接输出，
修复率为保守下界。

\subsection{工具链适配层}\label{sec:toolchain}

系统对 Yosys/SymbiYosys 工具链的解析失败采用“预防 + 降级”两层设计：其一，断言安全子集
（Gate-1 实测收敛）限定 immediate assert/assume、if 门控、打拍跨周期等结构，从源头避免并发
SVA（assert property、\$past、|->、\#\#n）等 Yosys/iverilog 不支持的结构进入形式验证；其二，
前置兼容性预检（scripts/check\_rtl\_compat.py）在 LLM 调用前扫描设计/断言文件中的不兼容结构
（并发 SVA、real 类型等），fail-fast 给出可操作改写建议，替代“运行中才崩溃”。验证阶段，
工具超时/解析错误由 evaluator 归类为 INCONCLUSIVE 并记录（不中断流水线），与 golden-first
规则共同构成降级路径；主协议中证明/k-归纳对门控断言返回 UNKNOWN 的问题即由此路径处理
（见 \S\ref{sec:judgment}）。

% ---------- 实验设置 ----------
\section{实验设置}\label{sec:setup}

\subsection{数据集 BugBench-PS}

BugBench-PS 包含 34 个 L3 样本（s04--s37），覆盖 6 个开源教学级模块：fifo\_sync、uart\_tx、
uart\_rx、axi\_lite\_slave、fsm\_ctrl 与 counter\_alu，错误类别共 7 类：状态迁移（24 样本）、
握手（12 样本）、复位（18 样本）、FIFO 满空（15 样本）、边界回绕（21 样本）、位宽截断
（9 样本）与边沿（3 样本），合计每设置 102 条定位记录。每个样本经 golden 双重校验（buggy
设计在 BMC 下 FAIL，golden 设计 PASS），并通过唯一性/结构审计（34/34 通过）。断言行、信号行
与随机行启发式基线全部失败（0\% 或 0.50\% 期望），证明注入缺陷位于功能/时序逻辑而非断言块。

\subsection{评测协议与成本记账}

主实验为干净口径下的 34 样本 $\times$ 3 设置（A/B/C）$\times$ 3 seeds $=$ 306 次真实 LLM
修复运行（DeepSeek v4-flash，C 102/102，temperature 0.2，OpenAI 兼容端点）。初版实验以
MiniMax M3 在四设置上运行 408 次，但审计发现三处 ground-truth 泄漏（见 \S\ref{sec:hygiene}），
其数字整体作废，仅作为卫生教训的对照引用；本文所有结论基于无泄漏重跑。修复成功由 BMC 加
golden 双重校验裁决。成本按 API token 记账（DeepSeek v4-flash 公开单价：输入 0.14 美元/
百万 token、输出 0.28 美元/百万 token，为估算值而非平台账单），干净主实验合计 1.65 美元
（A/B/C/D=408 次运行，平均 0.0041 美元/次，分设置 A 0.60/B 0.40/C 0.28/D 0.37 美元）。
L2 门禁对照使用 24 个 L2 样本 $\times$ A/B/C（72 次调用，见 \S\ref{sec:l2}）；D 设置（故障锥）
以同协议复跑完成（102 次调用，成本 \$0.37，定位 73.5\%，行号判据修复后）。运行采用并行调度：
8 并发、420 秒硬超时与断点续跑。

\subsection{可复现性}

全部实验基于开源工具链：iverilog 12.0、verilator 5.020、yosys 0.33、SymbiYosys（sby，
smtbmc+z3 引擎）与 Python 3.12，运行于 WSL。证据链生成命令为 \texttt{python3
scripts/build\_evidence.py --samples s04-s37 --real --window 8}；干净主实验（306 条）、L2 对照
（72 条）与 D 复跑（102 条）均经 \texttt{scripts/run\_experiments.py --provider deepseek}
执行，结果分别落盘到 \texttt{experiments/runs/leakfix\_0..7.json}、
\texttt{leakfix\_l2.json} 与 \texttt{leakfix\_D\_0..7.json}；充分性量化、T2 审计与 BMC 重验
分别由 \texttt{verify\_sufficiency.py}、\texttt{t2\_audit.py}、\texttt{reverify\_bmc.py} 完成。
所有 LLM 调用强制写入 token 账本（\texttt{experiments/runs/token\_ledger.jsonl}），成本按
输入 $\times$ 输入单价 + 输出 $\times$ 输出单价线性估算（DeepSeek v4-flash：0.14/0.28 美元
每百万 token，标注为估算口径）。数据集 BugBench-PS（34 样本，6 模块 $\times$ 7 错误类型）与
复现命令链一并开源，见仓库 \texttt{REPRO.md}。\textbf{数据溯源：}正文各表/图与数据文件的对应
关系为：表~\ref{tab:main}（主实验）← \texttt{leakfix\_0..7.json} +
\texttt{leakfix\_D\_0..7.json}；表~\ref{tab:hygiene_full}（卫生消融）←
\texttt{exp\_hygiene\_all.json}；表~\ref{tab:cross_model}（跨模型）←
\texttt{exp\_mm\_all\_full.json}（MiniMax）与 \texttt{exp\_hygiene\_all.json}（DeepSeek 卫生化）；
表~\ref{tab:error_matrix} ← \texttt{exp\_hygiene\_all.json} 聚合；深子集（第~\ref{sec:deep}~节）
← \texttt{exp\_cdeep\_ds\_*.json} 等。旧口径文件（如 \texttt{clean\_cross\_stats.json}）仅
保留作历史对照，不计入正文数字。

\subsection{非 LLM 基线}

为确立 LLM 定位的非平凡性，在相同精确匹配口径（候选行等于 buggy\_inject\_line，34 样本）下
评测三类简单启发式：首个断言行、任意断言行、首个提及失败信号的证据行，以及随机行期望。结果
全部为 0\%（随机行期望 0.50\%）；干净口径下 DeepSeek v4-flash 四设置的定位精度 60.8--87.3\%
远超这些基线，最高的 A（87.3\%）约为随机期望的 175 倍。注意：旧泄漏版曾以 MiniMax 的
47.1--61.8\% 作此比较，该数字已被卫生审计作废。

% ---------- 实验结果与分析 ----------
\section{实验结果与分析}\label{sec:results}

\subsection{结果概览}

\begin{table}[!t]
\centering
\caption{四设置干净主实验结果（DeepSeek v4-flash，n=102 每设置；BMC 回归通过率按有限深度
BMC 判据；卫生化后口径见第~\ref{sec:evidence-hygiene}~节）}
\label{tab:main}
\scriptsize
\setlength{\tabcolsep}{3pt}
\renewcommand{\arraystretch}{1.12}
\begin{tabular}{p{0.17\columnwidth}cp{0.11\columnwidth}p{0.11\columnwidth}p{0.12\columnwidth}p{0.10\columnwidth}p{0.09\columnwidth}}
\toprule
\textbf{设置} & \textbf{n} & \textbf{定位Top-1（卫生化前）} & \textbf{定位Top-1（卫生化后）} & \textbf{BMC通过率} & \textbf{Tokens} & \textbf{成本(USD)} \\
\midrule
A 原始日志 & 102 & 87.3\% & 95.1\%（97/102） & 100\% & 2.486M & 0.60 \\
B 结构化证据 & 102 & 75.5\% & 94.1\%（96/102） & 100\% & 1.672M & 0.40 \\
C 反例语义化 & 102 & 60.8\% & 92.2\%（94/102） & 100\% & 1.48M & 0.28 \\
D 故障锥 & 102 & 73.5\% & 95.1\%（97/102） & 100\% & 1.601M & 0.37 \\
\bottomrule
\end{tabular}
\end{table}

\begin{figure}[!htbp]
\centering
\includegraphics[width=0.98\columnwidth]{fig_setting_loc_cost.pdf}%
\caption{各设置的定位 Top-1 精度（柱，卫生化前）与 LLM 成本（折线，主实验口径）。四设置均为
无泄漏 DeepSeek v4-flash 运行；卫生化前 A vs C 显著（McNemar p=7.4e-6，Holm p=4.5e-5），
A vs D 亦显著（Holm p=0.047）；卫生化后（第~\ref{sec:evidence-hygiene}~节）四设置配对检验
全部不显著（A/C p=0.375、A/D p=1.0）。}
\label{fig:loc_cost}
\end{figure}

表~\ref{tab:main} 与图~\ref{fig:loc_cost} 汇总了四设置的主实验结果。\textbf{发现 1：}统一
BMC 判据下，无泄漏口径四种证据表征的 BMC 回归通过率均达 100\%——证据表征不决定\emph{能否}
修复，而决定\emph{多准、多便宜}；但卫生消融进一步表明，决定\emph{多准}的主要是证据呈现卫生
（行号基准、断言行语义），而非表征本身。

\textbf{发现 2：}卫生化前（原始呈现口径，DeepSeek v4-flash，无泄漏）定位精度为 A 87.3\%
（Wilson 95\% CI 79.4--92.4）$>$ B 75.5\%（66.3--82.8）$>$ 故障锥（D）73.5\%（64.2--81.1）
$>$ 反例语义化（C，DS 摘要版）60.8\%（102/102，CI 51.1--69.7）；该口径下 A vs C（p=7.4e-6，
Holm p=4.5e-5）与 A vs D（p=0.0094，Holm p=0.047）显著，其余四对 Holm 校正后均不显著。经
证据卫生消融（第~\ref{sec:evidence-hygiene}~节）后，四设置定位精度收窄为 A 95.1\%
（[89.0,97.9]）$=$ D 95.1\%（[89.0,97.9]）$>$ B 94.1\%（[87.8,97.3]）$>$ C 92.2\%
（[85.3,96.0]）（括号为 Wilson 95\% CI），极差从 29.4~pp 降至 2.9~pp，且配对检验全部不显著
（A/C p=0.375、A/D p=1.0、C/D p=0.45）——证据呈现卫生而非证据表征本身主导了卫生化前的差异。
成本（主实验口径）D 最低（0.37 美元/102 次），B 次之（0.40 美元），A 最高（0.60 美元）；
卫生化消融口径下为 D \$0.24、B \$0.29、C \$0.34、A \$0.37（见第~\ref{sec:evidence-hygiene}~节）。

\textbf{缺失样本分析：}C 设置缺失 7 个样本：s33（axi\_lite\_slave，状态跳转）全部 3 个 seeds、
s34（axi\_lite\_slave，边界回绕）seed 1、s37（uart\_rx，边界回绕）全部 3 个 seeds。缺失非随机：
全部为协议模块（axi/uart），其 buggy.v 更大（212--308 行），semantics.json 中的 fault\_cone
信号更多（14--31 个）。该 7 条已由 API 稳定后补跑协议：使用与原始运行相同的模型版本
（deepseek-v4-flash-0731）、温度（0.2）、反馈协议（F+hygiene）和 prompt 模板，补跑日期
2026-08-07。补跑数据与原始运行在统计分析中合并，合并前已确认补跑前后的统计结果无实质差异
（配对检验 p>0.1）。exp\_c\_ds\_full.json 现有完整 102/102。此前缺失源于 DeepSeek API 墙钟
超时（wall-clock timeout，非上下文容量限制——实测单次调用 input tokens 仅 7--14K，远低于
1M 上下文窗口），补跑后无样本缺口。\textbf{补跑合并的统计处理：}102 对配对样本中包含 7 对
补跑数据；配对 McNemar 以（样本，seed）为配对单元，补跑对与原始对同样独立计数；“无实质差异”
的判定为：仅对 7 条补跑样本按其原始缺失时点前后各运行集的定位结果做配对比较，未发现方向性
翻转（p>0.1）。由于缺失机制为 API 墙钟超时（与样本内容无关的传输层事件，非模型拒绝或上下文
溢出），我们认为不存在选择偏差；该 7 条的样本/seed 清单见仓库 REPRO.md 附录。\textbf{seed
独立性说明：}三组 seed 通过在 prompt 中注入不同的 seed 标识（占位符）并要求 temperature=0.2
的 API 采样获得；由于采样器存在内在随机性，相同 prompt 在不同请求间输出不完全相同，三组 seed
视为近似独立的重复抽样。配对统计按 102 个（样本，seed）单元独立计数；我们同时报告跨 seed 的
聚合定位率波动（见下文），作为独立性是否成立的间接证据。

\textbf{统计显著性（卫生化前口径）：}对逐样本定位结果按（样本，seed）配对做 McNemar 精确检验，
报告全部六对对比并经 Holm 校正。
A 对 C 显著（$p=7.4\times10^{-6}$，率差 $+26.5$~pp，95\%~CI $[+15.0,+37.9]$，Holm 校正后
$p_{\mathrm{holm}}=4.5\times10^{-5}$）——原始日志（A）的定位精度显著高于 DS 摘要版反例语义化
（C）；A 对 D 校正后亦显著（$p=0.0094$，$+13.7$~pp，$[+3.0,+24.5]$，$p_{\mathrm{holm}}=0.047$）。
经 Holm 校正后其余四对均不显著：
A 对 B $p=0.0357$（$+11.8$~pp，$[+1.2,+22.3]$，$p_{\mathrm{holm}}=0.110$）、
B 对 C $p=0.0275$（$+14.7$~pp，$[+2.1,+27.3]$，$p_{\mathrm{holm}}=0.110$）、
C 对 D $p=0.0533$（$-12.7$~pp，$[-25.5,+0.0]$，$p_{\mathrm{holm}}=0.110$）、
B 对 D $p=0.8450$（$+2.0$~pp，$[-10.0,+13.9]$，$p_{\mathrm{holm}}=0.845$）。
旧泄漏版报告的 B 对 A $p=0.0035$ 显著结论已被推翻。

\textbf{统计显著性（卫生化后口径）：}对卫生化后数据（exp\_hygiene\_all.json，A/B/C/D 各 102 对）
同样按（样本，seed）配对做精确 McNemar：A/C（不一致对 4/1）p=0.375、A/D（2/2）p=1.0、B/C
（2/0）p=0.5、C/D（2/5）p=0.45、A/B（2/1）p=1.0、B/D（2/3）p=1.0，\textbf{全部不显著}。
因此，卫生化后的四设置差异（极差 2.9~pp）不具统计意义，本文不宣称任何卫生化后表征排序的
显著性；该结果与第~\ref{sec:evidence-hygiene}~节的机制分析一致——卫生化前的显著差异主要源于
证据呈现层的行号/断言行语义缺陷。

seed 独立性方面，推理使用 temperature=0.2 且 prompt 携带 seed 标识（\S\ref{sec:setup}），
聚合定位率跨 seed 波动为：A 2.9~pp（88.2\%--85.3\%）、D 5.9~pp（76.5\%--70.6\%）、
C 8.8~pp（64.7\%--55.9\%）、B 20.6~pp（85.3\%--64.7\%），说明三-seed 配对提供了真实而非
名义的抽样方差，其中 B 设置的噪声最大。全部对比均报告，包括不显著项。

\textbf{定位距离（loc\_dev）：}除二值的 top-1 判据外，我们还报告 LLM 自报候选行与缺陷真实行
的距离中位数（\textit{loc\_dev}，基于 LLM 输出中显式给出的行号，而非 diff 应用位置）。按
修复后的真实缺陷行重算，主数据集四设置 loc\_dev 中位数均为 0 行（跨设置稳定），均值分别为
A 0.5、B 2.0、D 1.8、C 11.3——C 的均值偏高是因为语义化证据偶尔误导 LLM 输出相距甚远的行号
（拉高均值），但其中位数仍为 0 行。注意：loc\_dev 与 loc\_top1（diff 应用的精确匹配率）是
两个独立指标：loc\_dev 衡量 LLM 对缺陷位置的\textit{自我意识}的精确度（其输出中的行号），
loc\_top1 衡量修补工具对\textit{实际缺陷行}的修复成功率（diff 应用的 hunk 匹配）。loc\_dev
中位数在四设置间的稳定（均为 0）说明 LLM 在多数运行中精确命中缺陷行，其感知偏差不受证据
表示显著影响——此发现与卫生化前 top-1 的显著差异互补：前者度量自我意识、后者度量工具端
应用结果。

注意：C 设置 102/102 条已于 2026-08-07 全部补齐（原缺失 7 条因 API 墙钟超时而非上下文限制），
其 Wilson 95\% CI 为 51.1--69.7\%（62/102）；统计检验基于 102 条配对样本，正式结论以此口径
为准。

\subsection{错误类型分析}

\begin{figure}[!htbp]
\centering
\includegraphics[width=0.98\columnwidth]{fig_error_setting_heatmap.pdf}%
\caption{错误类别 $\times$ 设置的定位 Top-1 精度热力图。困难类别排序（状态迁移/握手低于位宽
截断/边沿）在全部设置中成立。}
\label{fig:error_heatmap}
\end{figure}

\begin{table}[!t]
\centering
\caption{错误类别 $\times$ 设置定位精度（每设置 n：状态迁移 24、握手 12、复位 18、FIFO 满空
15、边界回绕 21、位宽截断 9、边沿 3）}
\label{tab:error_matrix}
\scriptsize
\setlength{\tabcolsep}{4pt}
\renewcommand{\arraystretch}{1.12}
\begin{tabular}{p{0.24\columnwidth}cccc}
\toprule
\textbf{错误类别} & \textbf{A} & \textbf{B} & \textbf{C} & \textbf{D} \\
\midrule
状态迁移 & \textbf{100.0\%} & 83.3\% & 66.7\% & 87.5\% \\
握手 & \textbf{66.7\%} & 50.0\% & 16.7\% & 50.0\% \\
复位 & 77.8\% & \textbf{88.9\%} & 66.7\% & 61.1\% \\
FIFO 满空 & \textbf{73.3\%} & 66.7\% & 60.0\% & 60.0\% \\
边界回绕 & \textbf{95.2\%} & 85.7\% & 66.7\% & 81.0\% \\
位宽截断 & \textbf{100\%} & 55.6\% & 88.9\% & 88.9\% \\
边沿 & \textbf{100\%} & 66.7\% & 33.3\% & \textbf{100\%} \\
\bottomrule
\end{tabular}
\end{table}

图~\ref{fig:error_heatmap} 与表~\ref{tab:error_matrix} 给出按错误类别分解的定位精度。
\textbf{发现 3：}单点语义错误易定位（位宽截断 A 100\%、边沿 A 100\%），跨状态与握手错误难定位
（握手 45.8\%，n=48；状态迁移在新判据下大幅提升至 84.4\%，n=96），握手难度排序在全部设置中
成立。干净口径下，A（原始日志）在状态迁移/位宽/边界类上表现最强（状态迁移 100\%、位宽截断
100\%、边界回绕 95.2\%），B 仅在复位类上最高（88.9\%）——与旧泄漏版“B 在边界回绕/位宽截断
最优”的结论相反（握手/边沿类每设置 n$\le$12，小样本需谨慎解读）。

\subsection{修复可信度：充分性与 T2 审计}

\begin{table}[!t]
\centering
\caption{变异分析充分性（针对 golden 设计）。被杀的操作定义：变异体被杀是指 BMC 能检测到变异
（即变异后的设计在有限深度 BMC 下产生反例，反例不同于原始设计）。这一指标衡量的是断言集对
缺陷的敏感性，而非修复流水线本身的修复能力。}
\label{tab:sufficiency}
\scriptsize
\setlength{\tabcolsep}{2pt}
\renewcommand{\arraystretch}{1.12}
\begin{tabular}{p{0.24\columnwidth}p{0.20\columnwidth}p{0.13\columnwidth}p{0.12\columnwidth}}
\toprule
\textbf{变异家族} & \textbf{变异体} & \textbf{被杀} & \textbf{比率} \\
\midrule
强变异（比较器反转） & 400 & 354 & 88.5\% \\
常量变异（参数常量替换） & 484 & 396 & 81.8\% \\
删除变异（断言移除） & 分层抽样（每模块 2～4 个断言） & 全部 PASS & 空泛性控制 \\
\bottomrule
\end{tabular}
\end{table}

变异分析结果见表~\ref{tab:sufficiency}。T2 确定性审计对干净口径 408/408 条修复全部通过：
接口变更 0、断言篡改 0、证据闭环失败 0；loc\_dev 中位数 0（A/B/C 精确命中 74.5\%）。补丁
最小化由均值 2.26 行、最大值 12 行佐证（unified diff 增删行合计）。

\subsection{成本与性能}\label{sec:cost}

四设置干净主实验总成本 1.65 美元（408 次运行，平均 0.0040 美元/次）；分设置为 A 0.60、
B 0.40、C 0.28、D 0.37 美元，token 分别为 2.486M、1.672M、1.48M、1.601M。C 的成本来自语义化
摘要文本而非轨迹长度：主数据集反例中位 6 拍（VCD 时钟事件合计 295 拍、最长 41 拍），轨迹较短，
窗口压缩收益有限。验证性能方面，BMC 重验与深度敏感性抽查（见 \S\ref{sec:judgment}）按样本
并行调度（8 并发）；单样本 golden 验证耗时以 axi 约 88 秒、uart\_rx 约 153 秒为上限。

\subsection{L2 门禁对照：假阳性率（探索性）}\label{sec:l2}

为量化数据集纯度门禁（弱测试平台可观测）对系统能力的区分度，我们构造 24 个 L2 样本（单周期
功能错误，弱 tb FAIL + formal fail 双检通过，覆盖 axi/counter/fifo/fsm/uart\_rx/uart\_tx
六模块），并以与主实验相同的无泄漏协议用 DeepSeek v4-flash 在 A/B/C 三设置下各评测一次
（72 次调用，成本 \$0.48）。结果：L2 的 BMC 回归通过率为 91.7\%（66/72），与旧泄漏版数值相同
但已由无泄漏协议复现——该通过率是稳定现象而非泄漏伪影；它与 L3 基线（100\%）仍差 8.3 个
百分点——本实验设计无法分离“弱 tb 可观测性”与“缺陷跨周期复杂性”两个因素——不宜将 L2--L3
通过率差异直接归因于弱 tb 门禁。

未通过的 6 次调用全部集中在两个 L2 样本（l2\_uarttx\_01 起始位握手失效、l2\_uarttx\_02 状态
跳转至 STOP），三设置均为 INCONCLUSIVE，失败模式一致：候选补丁通过 BMC 与 golden 校验，但弱
tb 仿真回归仍未通过（formal=pass、sim 挂），是 uart\_tx 协议交互类缺陷难以定位修复的又一证据。
按设置看，干净口径下 L2 定位精度为 A 58.3\%、C 45.8\%、B 37.5\%（n=24，小样本），与 L3 上
A>C>B 的排序方向一致——与旧泄漏版“B 在 L2 最高（50.0\%）且与 L3 B 优势方向一致”的结论相反。
该对照实验说明：\textbf{证据表征的相对优势主要体现于跨周期（L3）缺陷}，对单周期可观测（L2）
缺陷区分度有限；论文将 L2 门禁的纯度代价如实报告：91.7\% 的通过率意味着弱 tb 门禁无法按修复
能力对单周期缺陷分层。

\subsection{判据翻案的方法论教训}\label{sec:judgment}

\textbf{发现 4：}原 prove/k-归纳判据误判 78 个正确修复（golden 本身在 prove 下 UNKNOWN）。对
A/B/C/D 408 次运行按 BMC 判据重验（A/B/C 306 次中 303 条含 diff 记录重验 + 3 条网络失败补跑，
D 102 条同协议复跑）全部通过，BMC 回归通过率由旧 prove 判据的 73.5\%（225/306）修正为 100\%。
判据强度抽查进一步将 34 个样本的修复在 2 倍 BMC 深度下重验（fifo/uart\_tx/fsm/counter
12$\to$24、axi 16$\to$32、uart\_rx 24$\to$48），34/34 全部通过，未发现深度敏感性翻转。这一
翻案说明：修复判据必须先经 golden 一致性校验，否则判据缺陷会系统性污染修复率、定位精度与
成本等全部下游指标。PreCex 将 BMC 判据与 golden-first 规则固化为评测协议的一部分，作为可复用
的方法论产出。

\subsection{数据卫生审计}\label{sec:hygiene}

C 证据文本改由 DeepSeek 重新生成（MiniMax 仅保留为独立评分者）。此切换引入了一项\textbf{实验
设计的混淆}：旧泄漏版实验中 C 证据由 MiniMax 生成时定位为 59.8\%（102 条），切换为 DeepSeek
生成并修复数据卫生后为 60.8\%（102/102，行号判据修复后口径）——两者同时改变了证据生成模型
与数据卫生状态，不可直接归因于任一变量。虽可解释为“DS 摘要更精简、定位信息更少”，但更本质
的教训是：语义化证据的质量\textit{由生成模型决定}，而非证据表示本身的固有属性——以不同模型
生成同一设置的证据再比较其定位，等效于同时改变证据内容和语义化能力两个变量却当作单一因子
报告。论文如实承认此混淆，当前 A/B/C/D 四设置数字（第~\ref{sec:results}~节与表~\ref{tab:main}）
全部基于 DS 生成的干净 C 证据，结论限定为“DS 摘要版 C 的定位低于 A、D、B”，不宣称“语义化”
这一表示在跨模型下的普遍劣势。

\subsection{证据卫生消融：全四设置扩展}\label{sec:evidence-hygiene}

审稿人指出，原始卫生消融仅覆盖 C 设置的 28 个已知 miss，未验证 A/B/D 是否受同类证据呈现问题
（行号基准、断言行语义）影响。为此，本节将卫生消融扩展至全四设置（A/B/C/D），仅修改证据/提示
呈现：设计代码加绝对行号前缀、输出格式规定 \texttt{line} 为文件绝对行号、C/D 证据中
\texttt{failed\_line} 改名为 \texttt{assert\_line} 并显式标注断言行不等于缺陷行。四设置均完整
覆盖 34 样本$\times$3 seeds，共 408 次 DeepSeek v4-flash 调用。结果见表~\ref{tab:hygiene_full}。

\begin{table}[t]
\centering
\caption{全四设置证据卫生消融结果（DeepSeek v4-flash，行号判据修正后口径）}\label{tab:hygiene_full}
\scriptsize
\setlength{\tabcolsep}{2pt}
\renewcommand{\arraystretch}{1.12}
\begin{tabular}{p{0.19\columnwidth}p{0.14\columnwidth}p{0.14\columnwidth}p{0.11\columnwidth}p{0.11\columnwidth}p{0.10\columnwidth}}
\toprule
设置 & 卫生化前 loc & 卫生化后 loc & $\Delta$ & McNemar $\chi^2$ & $p$ \\
\midrule
A（原始日志）    & 87.3\% (89/102) & 95.1\% (97/102)  & +7.8pp & 4.08  & 0.039 \\
B（结构化证据）  & 75.5\% (77/102) & 94.1\% (96/102) & +18.6pp & 17.05 & $<\!10^{-4}$ \\
C（语义化证据）  & 57.9\% (59/102) & 92.2\% (94/102) & +34.3pp & 26.69 & $<\!10^{-6}$ \\
D（静态故障锥）  & ---             & 95.1\% (97/102) & ---     & ---   & --- \\
\bottomrule
\end{tabular}
\smallskip

{\footnotesize 注：D 无双样本基线（原 D 仅 15 样本子集，不可配对 McNemar）。A 行 $p$ 为精确
McNemar 双侧检验。C 行卫生化前 57.9\% (59/102) 为旧审计口径（含行号判据修复前计数），与主实验
表~\ref{tab:main} 的行号判据修复后口径 60.8\% 相差 2.9~pp，二者均保留以对齐各自来源；后文统计
一律采用行号判据修复后口径（A 87.3\%、B 75.5\%、D 73.5\%、C 60.8\%）。卫生化后定位（行号判据
口径）：A 95.1\% $=$ D 95.1\% $>$ B 94.1\% $>$ C 92.2\%，极差仅 2.9~pp（卫生化前 29.4~pp）；
Wilson 95\% CI 分别为 A/D [89.0,97.9]、B [87.8,97.3]、C [85.3,96.0]。}
\end{table}

\textbf{主要发现：}卫生化对所有设置均有显著正向效果（逐设置卫生化前后配对 McNemar：C +34.3pp、
$p<10^{-6}$；B +18.6pp、$p<10^{-4}$；A +7.8pp、精确 McNemar $p=0.039$）。C 受益最大——28/40
的 miss 在卫生化后消失，表明原始 C 的低定位精度主要是证据呈现问题而非模型推理能力不足。卫生化
后四设置定位差距大幅收窄（极差从 29.4~pp 降至 2.9~pp），排名为 A 95.1\% $=$ D 95.1\% $>$
B 94.1\% $>$ C 92.2\%。需要强调：对卫生化后数据按（样本，seed）配对做精确 McNemar，四设置两两
对比\textbf{全部不显著}（A/C 不一致对 4/1，p=0.375；A/D 2/2，p=1.0；B/C 2/0，p=0.5；C/D 2/5，
p=0.45；A/B 2/1，p=1.0；B/D 2/3，p=1.0），因此卫生化后的排序与 2.9~pp 极差不具统计意义，本文
不据此宣称任何表征的卫生化后优势。所有四设置的 BMC 回归通过率均为 100\%。卫生化后 C 的 8 个
剩余 miss（s15 全部 3 seeds、s19 2 seeds、s29 全部 3 seeds）均属删除型缺陷，确认精确行号判据对
删除型缺陷系统性过严。\textbf{混淆声明：}本文的行号卫生为联合处理（行号前缀、failed\_line
$\rightarrow$assert\_line 改名、输出格式约束同时引入），未做单因素分离；因此卫生效应是这三项
呈现改动的联合效应，不单独归因于任一项。

\subsubsection{C 设置逐样本审计与机制验证}

C 设置 40 个定位 miss 的逐样本审计揭示，绝大多数失败并非模型推理能力不足，而是\textbf{证据
呈现层的两个可修复缺陷}：
(1) \textbf{行号计数基准错位}：设计代码在 prompt 中不带行号，LLM 需自行数行，倾向按 module
相对行计数。由于 buggy 文件头注释占 8 行，16 个 miss 的报行满足 $\mathrm{loc} = \mathrm{truth}
- 8$，且这 16 条的定位理由文本全部正确指认缺陷机制；
(2) \textbf{断言行语义误导}：反例语义化证据中的 failed\_line 字段是 sby 报告的\textbf{断言行}
而非缺陷行，12 个 miss 的报行恰好等于 failed\_line（如 s08 报 137 而真值为 63、s17 报 234 而
真值为 122--123）。

为验证这两类归因，本文实施证据卫生消融：对 28 条已知 miss 重跑修复，仅修改证据/提示呈现——
(a) 设计代码加绝对行号前缀并在输出格式中规定 line 必须为文件绝对行号；(b) C 证据 failed\_line
改名为 assert\_line 并显式标注“断言行不等于缺陷行”。结果如表~\ref{tab:hygiene_detail} 所示：
\textbf{28/28 全部精确定位（行偏差为 0），27/28 修复通过 BMC 回归}。该结果证明：C 设置 60.8\%
的定位精度低估了系统真实能力，证据呈现层的行号语义纠偏即可将 C 的定位外推至 95.1\%（97/102；
叠加 near1\_other 组 9 目标 7/9 精确、9/9 修复，±1 行容差口径 96.1\%）。

\begin{table}[t]
\centering
\caption{证据卫生消融结果（DeepSeek v4-flash，行号判据修复后口径）}
\label{tab:hygiene_detail}
\scriptsize
\renewcommand{\arraystretch}{1.12}
\begin{tabular}{lccc}
\toprule
失败聚类 & miss 数 & 卫生修复后定位 & 修复通过 \\
\midrule
eq\_sem（断言行误导） & 12 & \textbf{12/12}（dev=0） & 11/12 \\
off8（行号 -8 偏移） & 16 & \textbf{16/16}（dev=0） & 16/16 \\
\midrule
合计 & 28 & \textbf{28/28} & 27/28 \\
\bottomrule
\end{tabular}
\end{table}

\textbf{握手样本 s15 的 2$\times$2 分解}：在证据中人工补充“应翻转未翻转信号”（TraceAnalyzer
的 key\_signal\_diffs）与逐步推理指令均未改变其精确行命中（4 变体 $\times$ 3 seeds 全部
loc\_top1=False），但 12 次运行的理由文本全部正确指认“S\_IDLE 的 tx\_start 分支未同步拉低
txd”，且\textbf{11/12 的 diff 实际修到了根因插入点}（在 tx\_start 分支内补 \verb|txd <= 1'b0|，
即缺失语句的语义插入位置，行 55--60），修复通过 11/12，定位偏差 $\le 2$ 行占 9/12、$\le 6$ 行
占 12/12。该结论在删除型缺陷样本 s29（缺失写分支）上复现：3 个 seed 的 diff 全部修到根因插入点
且修复通过，但精确行号仅命中 1 个 seed。这揭示第三类判据卫生问题：\textbf{删除型缺陷（被删语句
行在 buggy 中不存在）下，精确行号判据系统性过严}，任何正确修复都不可能精确命中被删行；定位
评估应结合修复语义正确性与插入点容差口径，而非单点行命中。

\textbf{状态迁移子集的机制级验证：}为确认状态迁移类样本的定位瓶颈确为证据行号卫生而非模型能力，
选取此前最难的 5 个状态迁移样本（s07/s08/s09/s18/s36）以 hygiene（行号前缀 + 断言行语义纠偏）
与 handshake\_v2（叠加 TraceAnalyzer 静默信号与逐周期 golden/buggy 值序列）两变体各 3 seeds
在 A 设置下补跑共 30 次：hygiene 变体 \textbf{15/15 精确定位（loc\_dev=0）且修复通过 15/15}；
handshake\_v2 变体 14/15 精确（唯一 miss 为网络 IncompleteRead，非证据问题）。该结果与
eq\_sem/off8 消融相互印证：状态迁移类在行号卫生修正下可完全精确定位（主实验四设置 66.7--100\%，
表~\ref{tab:error_matrix}；早期行号基准错位口径下约 30\%），瓶颈是证据呈现层而非模型能力。

\textbf{结构性修复模式对照（负结果）：}为检验“最小 diff 约束压制状态迁移/握手类修复”的假说，
对 6 个样本（s07/s08/s09/s15/s18/s36）以 plain（最小 unified diff）与 structural（放开最小化
约束，允许新增中间状态、辅助寄存器与重写转移条件，同时禁止接口与断言变更）两模式各 3 seeds
跑完整形式验证协议（36 次调用）：plain 精确 top-1 13/18（72.2\%）、structural 12/18（66.7\%）；
±1 行容差 plain 13/18、structural 15/18；两模式修复通过率均为 18/18（100\%）、接口变更均为 0。
结论：结构性模式未带来精确 top-1 增益（有数据支撑的负结果），修复力持平、容差略高；s15（删除型）
两模式均 0/3 精确命中但 3/3 修复通过，与删除型判据过严的结论一致。该对照表明状态迁移/握手类
定位瓶颈不在最小化约束本身，而在证据行号卫生与删除型判据口径。

该消融为判据卫生（第~\ref{sec:judgment}~节）补充了证据卫生维度：定位失败需先排除呈现层假象
（行号基准、断言行语义、删除型缺陷判据），再归因于模型或证据内容。

\subsection{模块与错误类型的交互效应}

\textbf{发现 5：}定位难度并非均匀分布，而是由缺陷类型与模块共同决定（干净口径 408 次运行，
BMC 判据）：协议交互类缺陷（uart\_tx/uart\_rx 的握手）仍是相对最难定位类别（A/B/C/D 分别为
66.7\%/50.0\%/16.7\%/50.0\%），而状态迁移类在新判据下大幅提升（A 100\%、D 87.5\%、B 83.3\%、
C 66.7\%），数值/位宽/边界类缺陷整体可定位性高（fifo\_sync 位宽截断 A/C 100\%、counter\_alu
状态迁移 A/C 100\%、axi\_lite\_slave 边界回绕 A 95.2\%）。在干净口径下，这类“算术/位宽/边界”
增益属于原始日志（A）而非结构化证据（B）——与旧泄漏版“B 增益集中在数值/位宽/边界”的结论相反，
再次说明泄漏虚增了 B 的相对定位优势。对协议交互类缺陷，证据表示难以改善定位精度，但形式再验证
闭环仍能保证 BMC 回归通过率 100\%。D（因果图）整体定位 73.5\%，仅在 counter\_alu 上最高
（100\%）——确定性提取的因果图并未带来相对原始日志的增益。

\subsection{深时序子集：反例时长与证据表征的交互（探索性）}\label{sec:deep}

为检验证据表征效应是否随反例时长变化，我们对 5 个深时序样本（s38--s42，反例 35--75 拍，仓库
samples/deep）以四设置 $\times$ 3 seeds 完成补测（60 次调用，成本 \$0.25）。深子集语义化采用
失败时刻尾部回溯的自适应窗口（window\_eff，见 \S\ref{sec:evidence}），其余设置与主实验协议一致。

\textbf{探索性发现：}全部 60 次补测（5 样本 $\times$ 四设置 $\times$ 3 seeds）均达 BMC 回归通过
（修复率 100\%），其中 s38 在 C/D 设置上定位较低（C 2/3、D 1/3，该样本本身极难），但四设置修复
均成功。剔除 s38 后四样本（s39--s42，n=12/设置）定位排序为 C 100.0\% $>$ A 83.3\% $=$ D 83.3\%
$>$ B 66.7\%，与主数据集（A 87.3\% $>$ B 75.5\% $>$ D 73.5\% $>$ C 60.8\%（102/102））发生反转；
含 s38 的全量（n=15/设置）排序为 C 93.3\% $>$ A 86.7\% $>$ D 73.3\% $>$ B 66.7\%，反转方向一致。
该反转说明：原始日志（A）在短反例（中位 6 拍）上更有利，而语义化（C）在长反例上更有利——证据
表征的相对优势随反例时长而改变，而非固定属性。需诚实声明其混杂因素：深子集的 C 设置同时采用了
自适应窗口（主数据集为固定 window=8），因此语义化增益可能与窗口策略耦合；反例时长与窗口策略的
单独效应留待未来消融。

\subsection{证据可解释性的探索性评估}\label{sec:interp}

\textbf{探索性评估：}为量化 C/D 证据表示的可解释性差异，我们以 5 维 Likert 量表（完整性、
可读性、因果性、可操作性、可信度，1--5 分）对 10 个样本的 C（反例语义化）与 D（静态故障锥
（static fault cone））证据进行独立评分。\textbf{评分独立性保证：}被评的 C 证据文本由 DeepSeek
v4-flash 生成（CexSemantizer，summary\_meta 可查证），主实验定位与修复亦由 DeepSeek 执行；
评分为 MiniMax M3 单模型独立完成（20 次调用，总成本 \$0.09）——MiniMax 不参与证据生成、不参与
修复、不参与定位，为完全外部的评分者。该设计消除了“评分者与被评对象相同模型”的潜在偏倚。在此
口径下，MiniMax 对 C/D 证据的五维绝对评分如下：C 证据完整性 3.33、可读性 2.33、因果性 3.11、
可操作性 3.00、可信度 2.44；D 证据完整性 2.56、可读性 2.56、因果性 2.11、可操作性 2.22、
可信度 2.11。MiniMax 评分显示 C（语义化）在因果性（3.11 vs 2.11）和可操作性（3.00 vs 2.22）
上明显优于 D（因果图），而 D 的可读性略高（2.56 vs 2.33），总体 C 更具解释优势。该结果与行为
代理（定位精度 loc\_top1，见第~\ref{sec:results}~节）方向一致，但论文不将单模型主观评分解读为
客观可解释性度量，仅作为探索性参考。

\subsection{跨模型鲁棒性}

\begin{table}[!t]
\centering
\caption{跨模型对比（MiniMax M3 vs DeepSeek v4-flash；三-seed 严格配对；双方均在同一证据卫生
协议（行号前缀+断言行重命名，\S\ref{sec:hygiene}）下复跑，BMC 回归通过率按有限深度 BMC 判据；
三设置均完整覆盖 102/102）}
\label{tab:cross_model}
\scriptsize
\setlength{\tabcolsep}{1.5pt}
\renewcommand{\arraystretch}{1.12}
\begin{tabular}{p{0.13\columnwidth}p{0.16\columnwidth}p{0.16\columnwidth}p{0.12\columnwidth}p{0.12\columnwidth}}
\toprule
\textbf{设置} & \textbf{MM 定位} & \textbf{DS 定位} & \textbf{MM 修复} & \textbf{DS 修复} \\
\midrule
A 原始日志 & 92.2\% (94/102) & 95.1\%（97/102） & 99.0\% & 100\% \\
B 结构化证据 & 92.2\% (94/102) & 94.1\%（96/102） & 97.1\% & 100\% \\
C 反例语义化 & 94.1\% (96/102) & 92.2\%（94/102） & 100\% & 100\% \\
\bottomrule
\multicolumn{5}{p{0.9\columnwidth}}{\footnotesize 三设置均为 102 对完整配对；A 行配对 McNemar 基于完整 102 对。}
\end{tabular}
\end{table}

\textbf{发现 6：}定位精度的模型依赖性与卫生调节通过跨模型对比体现（表~\ref{tab:cross_model}）。
主实验在无泄漏口径下用 DeepSeek v4-flash（temperature=0.2）对 34 个 L3 样本的 A/B/C/D 四设置
各跑 3 个 seeds（408 次调用，成本 \$1.65，全部 BMC 三通过，T2 408/408）：卫生化前总体定位
74.3\%，按设置 A 87.3\% > B 75.5\% > D 73.5\% > C 60.8\%（102/102），A vs C 显著（Holm
p=4.5e-5），A vs D 亦显著（Holm p=0.047）。与旧泄漏版（A 51.0\%、B 54.9\%、C 62.7\%）相比，
134 处 loc\_top1 翻转（行号判据修复后口径）且“B 显著最优”被推翻——泄漏主要虚增了 B 的相对
定位优势。跨模型对比在同一证据卫生协议下完成：初版 MM 复跑带行号前缀与断言行重命名
（\S\ref{sec:hygiene}），而 DS 基线未卫生化，C 设置因此呈现 +33.3~pp 的伪差距；为消除该对等性
缺陷，本文改用 DS 卫生化复跑数据（第~\ref{sec:evidence-hygiene}~节）与 MM 严格配对。公平口径下
三设置均无显著模型差异（配对 McNemar，均 102 对完整配对，基于 exp\_mm\_all\_full.json 与
exp\_hygiene\_all.json 实测）：A 92.2\%（94/102）vs 95.1\%（97/102），不一致对（DS 优/MM 优）
4/1，p=0.375；B 92.2\%（94/102）vs 94.1\%（96/102），3/1，p=0.625；C 94.1\%（96/102）vs
92.2\%（94/102），1/3，p=0.625。BMC 回归通过率跨模型一致（100\%）。在卫生化对等条件下，未观察
到模型间显著差异；初版“MM 显著优于 DS”的交互效应不成立。这支持本文核心主张——证据表征的效应
应限定于所测模型与协议，论文不宣称跨模型普遍性。

\textbf{方法论含义：}跨模型对比必须先统一修复判据口径；若以旧 prove 判据（通过率 73.5\%）对比
BMC 判据（100\%），会错误得出“不同模型修复率差异显著”。统一 BMC 口径后该差异消失。这进一步
支持本文的核心主张：在判据可靠的前提下，证据表征的相对效应应限定于所测模型，论文不宣称跨模型
普遍性。

% ---------- 结论 ----------
\section{结\quad 论}

本文提出 PreCex，一个开源、反例驱动的综合前 RTL 缺陷定位与修复流水线，并系统研究了此前反例
理解类系统未受控的证据表征维度。在干净口径（无泄漏）下，以 DeepSeek v4-flash 在 34 样本跨周期
（L3）基准上完成 408 次主实验（A/B/C/D 四设置（C 102/102，其余各 102）$\times$ 3 seeds，均为
无泄漏口径）与 72 次 L2 门禁对照，主要结论如下：

\begin{itemize}
\item 统一 BMC 判据（golden-first）下，四种证据表征在实验筛选口径下的 BMC 回归通过率均达
100\%（有限深度口径）。卫生化前（原始呈现口径）定位精度 A 87.3\% > B 75.5\% > D 73.5\% >
C 60.8\%，A vs C（Holm p=4.5e-5）与 A vs D（Holm p=0.047）显著；仅修正行号基准与断言行语义
等呈现问题后（全四设置 408 次），定位精度收窄为 A 95.1\% $=$ D 95.1\% > B 94.1\% > C 92.2\%
（极差 2.9~pp），配对检验全部不显著（A/C p=0.375、A/D p=1.0）——证据呈现卫生而非证据表征
本身主导了卫生化前的显著差异。深时序子集（s39--s42）上排序反转为 C 100\% > A 83.3\% $=$
D 83.3\% > B 66.7\%（探索性，第~\ref{sec:deep}~节）。
\item 评测双卫生是可靠结论的前提：ground-truth 泄漏（三处）使“B 显著最优”成为伪结论（剥离
后 134 处定位翻转、结论反转）；prove/k-归纳判据失效会误判 78 个正确修复。数据卫生（ground-
truth 隔离）与判据卫生（golden-first）已固化为评测协议。
\item 修复可信度有据可依：强变异 88.5\%、常量变异 81.8\% 被杀，空泛性控制通过，T2 审计
408/408 通过（接口/断言/证据闭环零违规），定位偏差（loc\_dev）中位数 0、精确命中 74.5\%
（A/B/C）。
\item 跨模型鲁棒性：BMC 回归通过率跨模型一致（100\%）；在同一证据卫生协议下，MiniMax M3 与
DeepSeek v4-flash 的定位精度三设置均无显著差异（A 92.2\% vs 95.1\%、B 92.2\% vs 94.1\%、
C 94.1\% vs 92.2\%，配对 McNemar p=0.375/0.625/0.625），初版“MM 显著优于 DS”源于证据卫生
状态不对等而非模型能力——在卫生化对等条件下，定位差异由证据表征与卫生状态驱动，并非单一模型
或单一表征的固有属性。
\end{itemize}

未来工作包括：(1) 将动态自适应反例窗口（失败时刻尾部回溯）正式纳入主数据集评测，并对深时序
子集上的语义化增益做窗口策略与反例时长的正交消融（第~\ref{sec:deep}~节已提供初步证据：深子集
上 C 100\% > A 83.3\% $=$ D 83.3\% > B 66.7\%）；(2) 放宽最小化约束，支持结构性重写候选
（拆分状态、增加中间态）并与 delta-debugging 结合；(3) 顶层连接性审计：将修复放回原网表/SoC
交互环境验证模块外副作用；(4) 多候选并发验证与在线仲裁（探索与利用）；(5) 向工业协议与更大
断言集扩展、更强变异家族、多提供商多模型评测、GROVE 式跨样本知识层，以及从“生成最小代码改动”
向“生成更强的不变式约束”演进的修复范式；(6) 引入真正的根因诊断引擎——基于模型检验的反例
轨迹切片与基于 SAT 的极小反例核心提取——使形式验证从“裁判”升级为“诊断指导”，将修复范式从
试错式（trial-and-error）推进为诊断驱动（diagnosis-driven）。

\subsection{局限性}\label{sec:limitation}

\textbf{架构边界：}本方法的定位是“快速生成并验证候选补丁”（candidate patch generator），
而非“完全消除设计错误”。具体而言：(1) 修复成功定义为“反例消失”而非“不变式成立”——BMC 是
有限深度的穷举测试，不是完备性证明，深度上界以外的错误仍可能存在；(2) “行级定位 + 最小 unified
diff”的约束偏好局部修复——状态机拆分、增加中间态或重写状态转移条件等结构性重写可能需要十几行，
超出最小化约束，状态迁移/握手类在新判据下分别为 84.4\% 与 45.8\%（表~\ref{tab:error_matrix}）；
结构性修复对照（36 次，见 \S\ref{sec:evidence-hygiene}）显示放开最小化约束并未带来 top-1 定位
增益（plain 72.2\% vs structural 66.7\%），但两模式 BMC 回归通过率均仍为 100\%——BMC 浅深度下
小补丁可能绕过触发路径（弱修复），由 delta-debugging 与 T2 审计部分约束；(3) 验证为模块级，
无顶层互联/协议栈审计，修复在模块外的副作用未覆盖；(4) 单流水线单候选，3 seeds 是离线统计
而非在线多候选仲裁。这四点均列入未来工作。

\textbf{内部效度：}所有结果以 BMC 为主判据；对带门控时序断言的模块，BMC 探测到缺陷深度加 2
并与 golden 验证交叉核对，但 BMC 深度上界并非完备性保证（深度敏感性抽查：34 个样本的修复在
2 倍 BMC 深度下全部通过，见 \S\ref{sec:judgment}）。需要诚实声明：该抽查覆盖全部样本的 2 倍
深度而非更深边界，不能证明所有修复在更大深度下仍通过；BMC 深度的选择对绝对通过率有影响——本文
的结论是在 BMC 深度设为当前值的操作口径下得出的，报告的是该口径下的相对比较结果。判据翻案审计
量化并修正了假阳性（误判正确修复为失败），但假阴性（错误修复因深度不足被漏判）的风险未被系统性
量化，这是本判据的已知边界。LLM 推理是流水线唯一随机成分（3 seeds 配对统计），解析、提取与
验证均为确定性。修复成功定义要求编译、弱测试平台回归与 BMC 无反例；不改变断言集的补丁仍可能在
属性集之外改变行为，T2 审计通过接口与断言篡改检查约束该风险，但并非等价性检验。

\textbf{外部效度：}基准含 34 个教学级模块样本（fifo/uart/axi-lite/fsm/counter），不宣称工业
协议覆盖；样本为确定性错误注入生成而非真实缺陷分布；握手仍是最难定位类别（干净口径四设置
16.7--66.7\%），状态迁移在行号判据修复后显著改善（A/B/C/D 66.7--100\%），uart\_rx 的波特时序
常量变异杀死率弱（常量 40.9\%），是结构断言对时序常量漂移不敏感的已记录局限。常见模块可能出现在
LLM 训练语料中，本文以构造日期记录、逐样本来源与错误类别多样性缓解，但无法完全排除残余记忆。
干净主实验 408 次运行均使用 DeepSeek v4-flash（A/B/C/D 各 102）；MiniMax M3 在同一卫生协议下
复跑——A/B/C 各 102/102（92.2\%、92.2\%、94.1\%）——与 DS 卫生化数据（A 95.1\%、B 94.1\%、
C 92.2\%）配对后三设置均无显著模型差异（配对 McNemar p=0.375/0.625/0.625），故跨模型结论限定
于所测模型与干净协议，绝对值可能随模型版本变化。另，L2 门禁对照的干净结果见 \S\ref{sec:l2}。

\textbf{BMC 深度与修复充分性：}本文以有限深度 BMC 判据（操作定义：给定深度内无反例）判定修复
成功。深度 2 倍抽查（fifo/fsm/counter 12$\to$24、axi 16$\to$32、uart\_rx 24$\to$48）34/34
通过，未发现深度敏感性翻转，但未穷尽所有深度。我们不主张有限深度 BMC 等同于不变式证明——反例
在有限深度内消失可能是深度不足的假象，而非真正的功能正确。论文将“修复成功”严格限定为“有限深度
BMC 回归通过”这一操作定义，不做更强的完备性承诺。

\textbf{构造效度：}loc\_top1 记录模型首选中候选行是否落入注入缺陷的参考 diff；BMC 回归通过率
按有限深度 BMC 判据计量；成本按 API token 记账；统计声明基于逐样本（样本$\times$seed）配对
McNemar 精确检验，报告全部六对对比并附 Holm 校正、Wilson 比例 95\% 置信区间与率差效应量
（含不显著项）。

% ---------------------------------------------------------------------------
% 参考文献（GB/T 7714，条目 8pt）
% ---------------------------------------------------------------------------
\begin{thebibliography}{99}
\fontsize{8}{10}\selectfont
\setlength{\itemsep}{0pt}

\bibitem{fvdebug} BAI Y, HAMAD G B, HO C T, et al. FVDebug: an LLM-driven debugging assistant for automated root cause analysis of formal verification failures[EB/OL]. (2025)[2026-08-11]. https://arxiv.org/abs/2510.15906.

\bibitem{grove} BAI Y, REN H. Learning to debug: LLM-organized knowledge trees for solving RTL assertion failures[EB/OL]. (2025)[2026-08-11]. https://arxiv.org/abs/2511.17833.

\bibitem{openllmformal} TRAN H T. Open-source LLM-driven formal verification: a multi-agent pipeline for RTL repair[EB/OL]. (2026)[2026-08-11]. https://arxiv.org/abs/2607.28877.

\bibitem{formalrtl} LI K, LI M, WEN X, et al. FormalRTL: verified RTL synthesis at scale[EB/OL]. (2026)[2026-08-11]. https://formind.netlify.app/files/papers/FormalRTL.pdf.

\bibitem{chipagents} WHALECHIP. WhaleChip uses ChipAgents to compress root cause analysis into minutes[EB/OL]. (2026)[2026-08-11]. https://semiwiki.com/eda/chipagents-ai/371464-whalechip-uses-chipagents-to-compress-root-cause-analysis-into-minutes/.

\bibitem{veripilot} WANG Y, LIU C, ZHANG J, et al. VeriPilot: an LLM-powered Verilog debugging framework[EB/OL]. (2026)[2026-08-11]. https://arxiv.org/abs/2606.23759.

\bibitem{veridebug} WANG N, YAO B, ZHOU J, et al. VeriDebug: a unified LLM for Verilog debugging via contrastive embedding and guided correction[EB/OL]. (2025)[2026-08-11]. https://arxiv.org/abs/2504.19099.

\bibitem{rtlfixer} TSAI Y D, LIU M, REN H. RTLFixer: automatically fixing RTL syntax errors with large language models[EB/OL]. (2023)[2026-08-11]. https://arxiv.org/abs/2311.16543.

\bibitem{fixbench} LI S, FU W, ZHAO Y, et al. Fixbench-RTL: a comprehensive benchmark for evaluating LLMs on RTL debugging[C]//Proceedings of the IEEE AsianHOST. Piscataway: IEEE, 2025: 1-6.

\bibitem{hdlfixbench} HDL-FixBench: a repository-scale RTL repair benchmark[EB/OL]. (2026)[2026-08-11]. https://openreview.net/forum?id=omi6IH5N42.

\bibitem{assertllm} FANG W, LI M, LI M, et al. AssertLLM: generating and evaluating hardware verification assertions from design specifications via multi-LLMs[EB/OL]. (2024)[2026-08-11]. https://arxiv.org/abs/2402.00386.

\bibitem{assertllm2} WU Y, FANG W, WANG J, et al. AssertLLM2: a comprehensive LLM benchmark for assertion generation from design specifications[EB/OL]. (2026)[2026-08-11]. https://arxiv.org/abs/2605.27472.

\bibitem{assertgen} LYU H, WANG Y, DU Y, et al. AssertGen: enhancement of LLM-aided assertion generation through cross-layer signal bridging[EB/OL]. (2025)[2026-08-11]. https://arxiv.org/abs/2509.23674.

\bibitem{lasso} LASSO: safety property generation with vacuity checking and coverage feedback[C]//Proceedings of the International Conference on Machine Learning and CAD. 2025.

\bibitem{lintllm} FANG Z, CHEN R, YANG Z, et al. LintLLM: an open-source Verilog linting framework based on large language models[EB/OL]. (2025)[2026-08-11]. https://arxiv.org/abs/2502.10815.

\bibitem{rule2drc} KIM J, BYUN J, HWANG D, et al. Rule2DRC: benchmarking LLM agents for DRC script synthesis with execution-guided test generation[EB/OL]. (2026)[2026-08-11]. https://arxiv.org/abs/2605.15669.

\bibitem{posteda} LIU P, XU N, TANG J, et al. Bridging the last mile of circuit design: PostEDA-Bench, a hierarchical benchmark for PPA convergence[EB/OL]. (2026)[2026-08-11]. https://arxiv.org/abs/2605.06936.

\bibitem{openrtlset} WANG J, WAN L J, PINGALI S, et al. OpenRTLSet: a fully open-source dataset for large language model-based Verilog module design[EB/OL]. (2026)[2026-08-11]. https://arxiv.org/abs/2606.10285.

\bibitem{edaschema} SHRESTHA P, AVERSA A, SAVIDIS I. EDA-Schema-V2: a multimodal schema, open datasets, and benchmarks for machine learning in digital physical design[EB/OL]. (2026)[2026-08-11]. https://arxiv.org/abs/2605.06952.

\bibitem{chiplingo} LI L, YU X, NI J, et al. ChipLingo: a systematic training framework for large language models in EDA[EB/OL]. (2026)[2026-08-11]. https://arxiv.org/abs/2604.27415.

\bibitem{yosys} WOLF C. Yosys open synthesis suite[EB/OL]. (2013)[2026-08-11]. https://yosyshq.net/yosys/.

\bibitem{sby} WOLF C. SymbiYosys (sby): a front-end for Yosys-based formal verification flows[EB/OL]. (2018)[2026-08-11]. https://github.com/YosysHQ/SymbiYosys.

\bibitem{zh1} 王浩仁, 崔展齐, 岳雷, 等. 基于冗余覆盖信息约简的软件缺陷定位方法[J]. 电子学报, 2024, 52(1): 324-337.

\bibitem{zh2} 曹鹤玲, 姜淑娟. 基于Chameleon聚类分析的多错误定位方法[J]. 电子学报, 2017, 45(2): 394-400.

\bibitem{zh3} 王建峰, 魏长安, 盛云龙, 等. 基于错误交互集的组合测试软件故障定位方法[J]. 电子学报, 2014, 42(6): 1173-1178.

\bibitem{zh4} 陈翔, 鞠小林, 万志, 等. 基于程序频谱的动态缺陷定位方法研究[J]. 软件学报, 2015, 26(2): 390-412.

\bibitem{zh5} 姜淑娟, 张旭, 王荣存, 等. 基于路径分析和信息熵的错误定位方法[J]. 软件学报, 2021, 32(7): 2166-2182.

\end{thebibliography}

\end{document}
